\mysec{Lista 2}
\newcounter{dois}
\stepcounter{dois}
\myexe{\thedois}
\bigline
\textbf{Sejam $(X, \mcal F, \mu)$ um espaço mensurável e $A \in \mcal F$. 
Mostre que a aplicação $\lambda: \mcal F \rightarrow \xbR$ dada por 
$\lambda(E)\coloneq \mu(E \cap A)$ é uma medida.}
\begin{proof}[\textbf{Dem:}]
Ora, sabendo que $\mu$ é uma medida, temos 
$\lambda(\emptyset) = \mu(\emptyset \cap A) = \mu(\emptyset) = 0$. Mais ainda, 
visto que para todo $B \in \mathcal F$, $\mu(B) \ge 0$, temos que dado $E \in \mathcal F$,
$E \cap A \in \mcal F$, ou seja, 
$$\lambda(E) = \mu(E \cap A) \ge 0.$$ 
Para verificar a propriedade de $\sigma$-aditividade, considere ${E_n}$ uma família em 
$\mathcal F$, dois a dois disjuntos. Como consequência da definição de ${E_n}$, temos que 
$\{E_n \cap A\}$ também é dois a dois disjuntos. Com isso,
\begin{align*}
   \lambda\left(\bigcup\limits_{n \in \mathbb N} E_n\right) &= 
   \mu\left(\left(\bigcup\limits_{n \in \mathbb N} E_n\right)\cap A\right) = 
   \mu\left(\bigcup\limits_{n \in \mathbb N} (E_n \cap A)\right),\\
   &=\sum_{n \in \mathbb N}\mu(E_n \cap A)=\sum_{n\in\mbb N}\lambda(E_n).
\end{align*}
Desse modo, segue que $\lambda$ é uma medida, como queríamos.
\end{proof}
\biglinend

\newpage
\stepcounter{dois}
\myexe{\thedois}
\bigline
\textbf{Sejam $(X, \mcal F)$ um espaço mensurável $\mu_1, \cdots, \mu_n$ medidas em 
$(X, \mcal F)$ e $a_1, \cdots, a_n$ reais não negativos. Mostre que a aplicação 
$\lambda: \mcal F \rightarrow \xbR$ dada por 
\begin{align*}
  \lambda(E) = \sum_{k=1}^{n} a_k\mu_k(E)
\end{align*}
é uma medida em $(X, \mcal F)$.}
\begin{proof}[\textbf{Dem:}]
Visto que para todo $k \in \mbb N_{\le n}$, temos que $\mu_k$ é uma medida em $(X, \mathcal F)$,
então, para cada uma das medidas, segue que:
\begin{enumerate}
  \item $\mu_k(\emptyset) = 0$; 
  \item $\mu_k(E) \ge 0, \text{ para todo } E \in \mcal F$; 
  \item Vale a $\sigma$-aditividade para sequências de conjuntos disjuntos.
\end{enumerate}
Desse modo,
\begin{align*}
  \lambda(\emptyset) = \sum_{k=1}^{n} a_k\mu_k(\emptyset) \stackrel{(1)}{=} 
  \sum_{k=1}^{n} a_k \cdot0 = \sum_{k=1}^n 0 = 0. 
\end{align*}
Mais ainda, visto que $a_1, \dots, a_n$ é uma sequência de números não negativos,
  temos por $(2)$, que $a_k \mu_k \ge 0$ para todo $k \in \mbb N_{\le n}$. Ou seja 
\begin{align*}
  \lambda(E) = \sum_{k=1}^n a_k \mu_k(E) \ge 0, \text{ para todo } E \in \mcal F.
\end{align*}
  Por fim, dada uma sequência $\{E_j\}_{j \in \mbb N}$ em $\mcal F$ disjunto dois a dois,
  temos, 
  \begin{align*}
    \lambda\left(\bigcup\limits_{j\in \mbb N} E_j \right) &= 
    \sum_{k=1}^{n} \left[a_k \mu_k\left(\bigcup\limits_{j\in \mbb N} E_j\right)\right]
    \stackrel{(3)}{=}\sum_{k=1}^{n} \left[a_k \sum_{j \in \mbb N}\mu_k(E_j)\right]\\ 
    &= \sum_{k=1}^{n} \left[\sum_{j \in \mbb N}a_k\mu_k(E_j)\right]
    \end{align*}
Como a soma em $k$ é finita, podemos trocar a ordem das somas, obtendo
\begin{align*}
    &= \sum_{j \in \mbb N} \left[\sum_{k \in \mbb N_{\le n}}a_k\mu_k(E_j)\right]
    = \sum_{j \in \mbb N} \lambda(E_j).
  \end{align*}
  Com isso, temos que $\lambda$ é uma medida em $(X, \mcal F)$.
\end{proof}

\biglinend
\newpage 
\stepcounter{dois}
\myexe{\thedois}
\bigline 
\textbf{
  Sejam $(X, \mcal F)$ um espaço mensurável e $(\mu_n)_{n \in \mbb N}$ uma sequência 
  de medidas em $(X, \mcal F)$ tal que $\mu_k(X) = 1$, para todo $k \in \mbb N$. 
  Mostre que a aplicação $\lambda: \mcal F \rightarrow \xbR$ dada por 
\begin{align*}
  \lambda(E) = \sum_{k \in \mbb N} \frac{\mu_k(E)}{2^k},
\end{align*}
é uma medida e $\lambda(X) = 1$.
}
\begin{proof}[\textbf{Dem:}] Vejamos que $\lambda(\emptyset) = 0$. 
Ora, $\lambda(\emptyset) = \sum_{k \in \mbb N} \frac{\mu_k(\emptyset)}{2^k}$. Já que 
  $\mu_k(\emptyset) = 0$ para todo $k \in \mbb N$, temos $\lambda(\emptyset) = 0$.
  Daqui, visto que para cada $E \in \mcal F$ e $k \in \mbb N$ temos $\mu_k(E)/2^k \ge 0$, 
então, para todo $E \in \mcal F$,
  \begin{align*}
    \lambda(E) = \sum_{k \in \mbb N} \frac{\mu_k(E)}{2^k} \ge 0.
  \end{align*}
  Com isso, vejamos que $\lambda$ é $\sigma$-aditiva. Dada uma sequência 
  $\{E_j\}_{j \in \mbb N}$ em $\mcal F$ disjunto dois a dois, temos, 
  \begin{align*}
    \lambda\left(\bigcup E_j\right) = 
    \sum_{k \in \mbb N} \frac{\mu_k \left(\bigcup E_j\right)}{2^k}.
  \end{align*}
  Já que vale a $\sigma$-aditividade para todo $\mu_k$, temos,
  \begin{align*}
    \lambda\left(\bigcup E_j\right) = 
    \sum_{k \in \mbb N}\left(\sum_{j \in \mbb N} \frac{\mu_k (E_j)}{2^k} \right).
  \end{align*}
  Como $\mu_k(E_j)/2^k \ge 0$ para cada $j,k \in \mbb N$,
 \begin{align*}
    \lambda\left(\bigcup E_j\right) = 
    \sum_{j \in \mbb N}\left(\sum_{k \in \mbb N} \frac{\mu_k (E_j)}{2^k} \right) = 
    \sum_{j \in \mbb N}\lambda(E_j).
  \end{align*}
 Por fim, visto que $\mu_k(X) = 1$ para todo $k \in \mbb N$, 
\begin{align*}
  \lambda(X) = \sum_{k \in \mbb N} \frac{\mu_k(X)}{2^k} 
  = \sum_{k \in \mbb N} \frac{1}{2^k} = 1. 
\end{align*}
Com isso, temos o que queríamos.
\end{proof}

\biglinend
\newpage 
\stepcounter{dois}
\myexe{\thedois}
\bigline
\textbf{Seja $(\mbb N, \mcal P(\mbb N))$ e $(a_n)_{n \in \mbb N}$ uma sequência 
de números reais estendidos não negativos. Defina $\mu: \mcal P(\mbb N) \rightarrow \xbR$
por $\mu(\emptyset) = 0$ e, para $E \neq \emptyset$,
\begin{align*}
  \mu(E) = \sum_{n \in E} a_n.
\end{align*}
Mostre que $\mu$ é uma medida em $\mbb N$. Reciprocamente, mostre que toda medida 
em $(\mbb N, \mcal P(\mbb N))$ é obtida desta maneira, para alguma sequência 
$(a_n)$ em $\xbR$.
}
\begin{proof}[\textbf{Dem:}]
  Pela definição de $\mu$, temos $\mu(\emptyset) = 0$. Mais ainda, como $a_n \ge 0$ para todo $n \in \mbb N$, dado 
  $E \in \mcal P(\mbb N)$, temos
  $\mu(E) = \sum_{n \in E} a_n \ge 0$. Daqui, vejamos que $\mu$ é $\sigma$-aditiva. Dada uma sequência 
  $\{E_n\}_{n \in \mbb N}$ em $\mcal P(\mbb N)$ disjuto dois a dois. Se $\bigcup E_j = \emptyset$, então $E_j = \emptyset$ para todo 
  $j \in \mbb N$. Portanto, $\mu(\bigcup E_j) = 0 = \sum \mu(E_j)$. Assumimos então, que $\bigcup E_j \neq \emptyset$, 
  desse modo,
  \begin{align*}
    \mu\left(\bigcup\limits_{j \in \mbb N} E_j\right) &= \sum_{n \in \bigcup E_j} a_n.
  \end{align*}
  Considere $J = \{n \in \mbb N\st  E_n = \emptyset\}$. Visto que $\{E_j\}$ é disjunto dois a dois,
  para todo $n \in \bigcup E_j$ existe um único $k \in \mathbb N$ tal que 
  $n \in E_k$, e, como para cada $n \in \bigcup E_j$, temos $a_n \ge 0$, segue que,
  \begin{align*}
    &= \sum_{n \in \bigcup E_j} a_n = \sum_{j \in \mbb N \setminus J} \left[\sum_{n \in E_j} a_n\right]
    = \sum_{j \in \mbb N \setminus J} \mu(E_j).
  \end{align*}
  Da definição de $J$, 
\begin{align*}
  &= \sum_{j \in \mbb N \setminus J} \mu(E_j) + \sum_{l \in J}\mu(E_l) = \sum_{j \in \mbb N} \mu(E_j).
\end{align*}
  Com isso, segue que $\mu$ é $\sigma$-aditiva. Já que $\mcal P (\mbb N)$ é uma $\sigma$-álgebra, $\mu$ é medida em 
  $(\mbb N, \mcal P (\mbb N))$.

  Agora, vejamos que toda medida em $\mbb N$ pode ser escrita como $\mu$. Considere $\lambda$ 
  uma medida em $(\mbb N, \mcal P (\mbb N))$. Seja $E \in \mcal P (\mbb N)$ tal que $E \neq \emptyset$.
  Então, 
  \begin{align*}
    \lambda(E) = \lambda\left(\bigcup\limits_{n \in E} \{n\}\right).
  \end{align*}
  Como $E \subseteq \mbb N$, $E$ é enumerável. Mais ainda, $\{n\}_{n \in E}$ é disjunto dois a dois, portanto, 
  da $\sigma$-aditividade, 
  \begin{align*}
    &= \lambda\left(\bigcup\limits_{n \in E} \{n\}\right) = \sum_{n \in E} \lambda(\{n\}),
  \end{align*}
  definindo $a_n = \lambda(\{n\})$, temos que cada $a_n \in \xbR$ é não-negativo, 
  como queríamos.
\end{proof}

\biglinend
\newpage 
\stepcounter{dois}
\myexe{\thedois}
\bigline
\textbf{Seja $X$ um conjunto não-enumerável. Considere a coleção 
$$\mcal F = \{A \subseteq X\st  A \text{ ou } A^C \text{ é enumerável}\}.$$ Defina $\mu: \mcal F\rightarrow \mbb R$
por 
\begin{align*}
  \mu(E) = \begin{cases}
    0, &\text{ se $E$ é enumerável;}\\
    1, &\text{ se $E^C$ é enumerável.}
  \end{cases}
\end{align*}
Mostre que $\mu$ é uma medida em $(X, \mcal F)$.
}
\begin{proof}[\textbf{Dem: }]
  Primeiro, note que $\mcal F$ de fato é uma $\sigma$-álgebra. E, visto 
  que $\emptyset$ é enumerável, $\mu(\emptyset) = 0$, e como para todo $E \in \mcal F$, temos $\mu(E) = 0$ 
  ou $\mu(E) = 1$, segue que $\mu(E) \ge 0$. Daqui, vejamos que $\mu$ é $\sigma$-aditiva. 
  Considere a sequência $\{E_j\}_{j \in \mbb N}$ em $\mcal F$ disjunto dois a dois. Se $\bigcup E_j$ é enumerável, 
  então $E_j$ é enumerável para todo $j \in \mbb N$, desse modo, 
  \begin{align*}
    \mu\left(\bigcup\limits_{j \in \mbb N} E_j\right) = 0 = \sum_{j \in \mbb N} \mu(E_j)
  \end{align*}
  Caso contrário, se $\bigcup E_j$ é não-enumerável, note que existe exatamente um $k \in \mbb N$ tal que $E_k$ é 
  não-enumerável, pois, supondo por contradição que existe $m \in \mbb N$ tal que $E_m$ é não-enumerável e $k \neq m$. 
  Como $E_k$ é não-enumerável e $E_k \in \mcal F$, então $E_k^C$ é enumerável. Entretanto, já que a sequência $\{E_j\}$
  é disjunto dois a dois, $E_m \subset E_k^C$, ou seja, $E_m$ deve ser enumerável, uma contradição. Com isso, garantimos 
  que existe exatamente um $k \in \mbb N$ tal que $E_k$ é não-enumerável. Daqui, como $\bigcup E_j \in \mcal F$ é não enumerável,
  pela definição de $\mcal F$, segue que 
  $\left(\bigcup E_j\right)^C$ é enumerável. Desse modo, visto que $\mu(E_j) = 0$ para todo 
  inteiro positivo $j \neq k$, $\mu(E_k) = 1$, temos
  \begin{align*}
    \mu\left(\bigcup\limits_{j \in \mbb N} E_j\right) = 1 = \mu(E_k) + \sum_{j \neq k} \mu(E_j) 
    = \sum_{j \in \mbb N} \mu(E_j).
  \end{align*}
  Ou seja, $\mu$ é $\sigma$-aditiva. Portanto, $\mu$ é uma medida em $(X, \mcal F)$.
\end{proof}


\biglinend
\newpage 
\stepcounter{dois}
\myexe{\thedois}
\bigline
\textbf{Seja $X$ um conjunto não-enumerável e considere o espaço mensurável $(X, \mcal P(X))$. 
Mostre que a aplicação $\mu: \mcal P(X) \rightarrow \xbR$ definida por 
\begin{align*}
  \mu(E) = \begin{cases}
    0, &\text{se $E$ é enumerável;}\\
    +\infty, &\text{se $E$ é não-enumerável.}
  \end{cases}
\end{align*}
Mostre que $\mu$ é uma medida em $(X, \mcal F)$.}
\begin{proof}[\textbf{Dem:}]
Visto que $\emptyset$ é enumerável, $\mu(\emptyset) = 0$. Já que para todo $A \in \mcal P(X)$, 
  $A$ é enumerável ou não-enumerável, temos que $\mu(A) = 0$ ou $\mu(A) = +\infty$, ou seja, 
  $\mu(A) \ge 0$. Daqui, vejamos que $\mu$ é $\sigma$-aditiva. Considere uma sequência 
  $\{E_j\}_{j \in \mbb N}$ em $\mcal F$ disjunto dois a dois. Se $\bigcup E_j$ é enumerável, 
  para todo $k \in \mbb N$, $E_k$ é enumerável. Sendo assim, 
  \begin{align*}
    \mu\left(\bigcup\limits_{j \in \mbb N} E_j\right) = 0 
    = \sum_{j \in \mbb N} 0 = \sum_{j \in \mbb N} \mu(E_j).
  \end{align*}
  Caso contrário, se $\bigcup E_j$ é não-enumerável, existe ao menos um $k \in \mbb N$ 
  tal que $E_k$ é não-enumerável. Portanto, 
  \begin{align*}
    \mu\left(\bigcup_{j \in \mbb N} E_j\right) = +\infty = \mu(E_k)
    = \sum_{j \in \mbb N} \mu(E_j).
  \end{align*}
  Com isso, segue que $\mu$ é $\sigma$-aditiva. Mais ainda, $\mu$ é uma medida.
\end{proof}


\biglinend
\newpage 
\stepcounter{dois}
\myexe{\thedois}
\bigline
\textbf{Considere no espaço 
mensurável $(\mbb N, \mcal P(\mbb N))$ a aplicação $\mu: \mcal 
P(\mbb N) \rightarrow \xbR$
definida por 
\begin{align*}
  \mu(E) = \begin{cases}
    0, &\text{ se $E$ é finito;}\\
    +\infty, &\text{ se $E$ é infinito.}
  \end{cases}
\end{align*}
A aplicaçã $\mu$ é uma medida?
}
\begin{proof}[\textbf{Dem:}]
  Visto que $\emptyset \in \mcal P(\mbb N)$ é finito, 
  $\mu(\emptyset) = 0$, mais ainda, como para todo 
  $E \in \mcal P (\mbb N)$, $E$ é finito ou infinito, segue que $\mu(E) \ge 0$. Com isso
  vejamos que $\mu$ obedece a $\sigma$-aditividade. Ora, considere uma sequência  
  $\{E_j\}_{j \in \mbb N}$ em $\mcal P(\mbb N)$ disjunto dois a dois. Se $\bigcup E_j$
  é finito, temos que $E_j$ é finito para todo $j \in \mbb N$, portanto, 
  \begin{align*}
    \mu\left(\bigcup_{j \in \mbb N} E_j\right) = 0 = \sum_{j \in \mbb N} 0
    = \sum_{j \in \mbb N} \mu(E_j). 
  \end{align*}
  Por outro lado, supondo que $\bigcup E_j$ é infinito, então, existe pelo 
  menos um $k \in \mbb N$ tal que $E_k$ é infinito. Desse modo, 
  visto que $\mu(E) \ge 0$ para todo $E \in \mcal P(\mbb N)$,
  \begin{align*}
    \mu\left(\bigcup_{j \in \mbb N} E_j\right) = +\infty = \mu(E_k)
    = \sum_{j \in \mbb N} \mu(E_j).
  \end{align*}
  Com isso, temos que $\mu$ é uma medida.
\end{proof}

\biglinend
\newpage 
\stepcounter{dois}
\myexe{\thedois}
\bigline 
\textbf{Sejam $(X, \mcal F, \mu)$ um espaço de medida e $(A_n)_{n \in \mbb N}$
uma sequência em $\mcal F$.}
\begin{enumerate}[label=(\alph*), wide, labelindent=0pt]
  \item \textbf{Mostre que o conjunto 
    $$\lim \sup A_n \coloneq \bigcap_{m=1}^\infty\left
    [\bigcup\limits_{n = m}^\infty A_n\right]$$ 
    consiste dos elementos de $X$ que pertencem a infinitos conjuntos da sequência $A_n$.}
    \begin{proof}[\textbf{Dem:}] Para verificar a proposição, dividiremos a demonstração 
      em duas partes. \\
      ($\implies$) Suponha $x \in \lim\sup A_n$, vejamos que $x$ pertence a infinitos 
      conjuntos da sequência $A_n$. Da hipótese, temos $x \in \bigcap_{m=1}^\infty 
      \left[\bigcup_{n = m}^\infty A_n\right]$, ou seja, $x \in \bigcup_{n=m}^\infty A_n$
      para todo $m \in \mbb N$. portanto, para todo $m \in \mbb N$, existe $n \ge m$ tal que 
      $x \in A_n$, i.e, $x$ pertence a infinitos conjuntos de $\{A_n\}$.\\
      ($\impliedby$) Suponhamos agora que $x$ pertence a infinitos conjuntos da sequência 
      $A_n$ e vejamos que $x \in \lim\sup A_n$. Ora, considere a sequência 
      $\{B_n\}_{n \in \mbb N}$ tal que $B_m = \bigcup_{n=m}^\infty A_n$. Da hipótese, 
      segue que $x \in B_n$ para todo $n \in \mbb N$, ou seja, 
      $x \in \bigcap\limits_{m=1}^\infty B_m = \lim\sup A_n$. 
    \end{proof}
  \item \textbf{Se $\mu\left(\bigcup\limits_{n \in \mbb N} A_n\right) < +\infty$, mostre que 
    $$\lim \sup \mu (A_n) \le \mu(\lim\sup A_n).$$
  Dê um exemplo em que a desigualdade é falsa quando $\mu\left(\bigcup
    \limits_{n \in \mbb N} A_n\right) = +\infty$.}
    \begin{proof}[\textbf{Dem:}] Definindo $F_m = \bigcup\limits_{n = m}^\infty A_n$, temos 
      $F_{m + 1} \subseteq F_{m}$, ou seja, $\{F_m\}$ é decrescente e mais ainda, 
        $\mu(F_1) = \mu\left(\bigcup\limits_{n \in \mbb N} A_n\right) < +\infty$. 
        Portanto, da continuidade da medida, 
        $$\mu\left(\bigcap_{m \in \mbb N} F_m\right) =
        \mu\left(\bigcap_{m=1}^\infty\left[\bigcup\limits_{n = m}^\infty A_n\right]\right)
        = \mu\left(\lim\sup A_n\right) = \lim_{m \rightarrow +\infty} \mu(F_m).$$
      Visto que cada $A_m \subseteq F_m$, temos $\mu(F_m) \ge \mu(A_m)$. Portanto, 
      \begin{align*}
        \mu(\lim\sup A_n) = \lim_{n \rightarrow +\infty}\mu(F_n) = 
        \lim\sup \mu(F_n) \ge \lim\sup \mu(A_n).
      \end{align*}
      Com isso, temos o que queríamos. Como exemplo em que a desigualdade é falsa, 
      considere $A_n = (n, +\infty)$. Já que $\mu(A_n) = +\infty$ para todo $n \in \mbb N$,
      então, $\lim\sup \mu(A_n) = +\infty$. Entretanto, visto que para todo $x \in \mbb R$,
      existe $M \in \mbb N$ tal que $x < M$, segue que $\lim\sup A_n = \emptyset$, ou seja, 
      $\mu(\lim\sup A_n) = 0$.
    \end{proof}
  \item \textbf{Mostre que o conjunto 
    $$\lim \inf A_n \coloneq \bigcup\limits_{m=1}^\infty\left[\bigcap\limits_{n=m}^\infty 
    A_n\right]$$
    consiste dos elementos de $X$ que não pertencem a apenas finitos conjuntos da  
    sequência $A_n$.}
    \begin{proof}[\textbf{Dem:}] Para verificar a proposição, dividiremos a demonstração 
      em duas partes.\\
      ($\implies$) Suponhamos que $x \in \lim\inf A_n$, vejamos que $x$ não pertence a apenas 
      finitos conjuntos da sequência $A_n$, i.e, existe $m \in \mbb N$ tal que para 
      inteiros $n \ge m$, temos $x \in A_n$. Ora, já que
      $$x \in \lim\inf A_n = \bigcup\limits_{m=1}^\infty\left[\bigcap\limits_{n=m}^\infty 
    A_n\right],$$
      segue que existe $m \in \mbb N$ tal que para inteiros $n \ge m$, temos $x \in A_n$,
      como queríamos.\\
      ($\impliedby$) Suponhamos que $x$ não pertence a apenas finitos conjuntos da sequência 
      $\{A_n\}$, vejamos que $x \in \lim\inf A_n$. Ora, da hipótese, temos que existe 
      $m \in \mbb N$ tal que para inteiros $n \ge m$ temos $x \in A_n$. Desse modo, 
      $$x \in \bigcup\limits_{m \in \mbb N}\left[\bigcap\limits_{n=m}^\infty A_n\right]
      = \lim \inf A_n$$
      como queríamos.
    \end{proof}
  \item \textbf{Mostre que $\mu(\lim \inf A_n) \le \lim\inf \mu(A_n)$.}
    \begin{proof}[\textbf{Dem:}] Definindo $F_m = \bigcap\limits_{n = m}^\infty A_n$,
      note que $F_m \subset F_{m+1}$, sendo assim, $\{F_m\}_{m \in \mbb N}$ é crescente. 
      Da continuidade da medida, segue que 
      \begin{align*}
        \mu\left(\bigcup_{m \in \mbb N} F_m\right)
        = \mu(\lim\inf A_n)
        = \lim_{m \rightarrow +\infty} \mu(F_m) = \lim\inf \mu(F_n).
      \end{align*}
      Já que $F_n \subset A_n$, temos $\mu(F_n) \le \mu(A_n)$, ou seja,
      \begin{align*}
        \mu(\lim\inf A_n) \le \lim\inf \mu(A_n),
      \end{align*}
      como queríamos.
    \end{proof}
\end{enumerate}
\biglinend

\newpage 
\stepcounter{dois}
\myexe{\thedois}
\bigline 
\textbf{Prove o \textit{Lema de Borel-Cantelli}: Sejam $(X, \mcal F, \mu)$ um espaço 
de medida e $(E_n)_{n \in \mbb N}$ uma sequência em $\mcal F$ tal que 
$\sum_{n \in \mbb N} \mu(E_n) < +\infty$. Então quase-$\mu$ todo ponto $x \in X$ pertence a, 
no máximo, um número finito de conjuntos da sequência $(E_n)_{n \in \mbb N}$}
\begin{proof}[\textbf{Dem:}] Note que, $x\in X$ pertence a, no máximo, um número 
  finito de conjuntos da sequência $(E_n)_{n \in \mbb N}$ $\mu$-q.t.p se 
  $\mu(\lim\sup E_n) = 0$. Daqui, visto que 
  \begin{align*}
    \mu\left(\bigcup\limits_{n \in \mbb N}
    E_n\right) \le \sum_{n \in \mbb N} \mu(E_n) < +\infty,
  \end{align*}
  definindo $F_m = \bigcup_{n = m}^\infty E_n$, 
  segue que $\{F_m\}$ é decrescente, ou seja, 
  pela continuidade da medida que $\lim \mu (F_n) = \mu(\bigcap F_m) = \mu(\lim\sup E_n)$. 
  Mais ainda, definindo $S_m = \sum_{n = m}^\infty \mu(E_n)$, como $\sum \mu(E_n) < +\infty$,
  segue que $\lim S_m = 0$. Portanto, da subaditividade, 
  \begin{align*}
    \mu(F_n) = \mu\left(\bigcup\limits_{n = m}^\infty E_n\right)
    \le \sum_{n = m}^\infty \mu(E_n) = S_m,
  \end{align*}
  ou seja, $\lim S_m = \lim \mu\left(\bigcup\limits_{n = m}^\infty E_n\right) 
  = \mu(\lim\sup E_n) = 0$.
\end{proof}
\biglinend

\newpage 
\stepcounter{dois}
\myexe{\thedois}
\bigline 
\textbf{Seja $(X, \mcal F, \mu)$ um espaço de medida finito.} 
\begin{enumerate}[label=(\alph*), wide, labelindent=0pt]
  \item \textbf{Mostre que se $E, F \in \mcal F$ e $\mu(E \Delta F) = 0$, 
    então $\mu(E) = \mu(F)$.}
    \begin{proof}[\textbf{Dem:}] Ora, visto que 
      $E \Delta F = (E \setminus F) \cup (F \setminus E)$ e 
      $(E\setminus F)\cap (F \setminus E) = \emptyset$, entâo, 
      \begin{align*}
        0 = \mu(E\Delta F) = \mu((E \setminus F) \cup (F \setminus E))
        = \mu(E\setminus F) + \mu(F\setminus E).
      \end{align*}
      Já que o espaço de medida é finito, temos $\mu(X)<+\infty$, ou seja $\mu(A) < +\infty$
      para todo $A \in \mcal F$, $(1)$. Com isso, 
      \begin{align*}
        \mu(E\setminus F) = -\mu(F\setminus E)
      \end{align*}
      Note que, se $\mu(F\setminus E) \neq 0$, teremos $\mu(F \setminus E) > 0$, i.e, 
      $\mu(E \setminus F) < 0$, uma contradição. Sendo assim, $\mu(F \setminus E) = 0$. 
      Com isso, já que 
      \begin{align*}
        F\setminus E = F \cap E^C = (F \cap E^C) \cup (F \cap F^C) = F\cap(E \cap F)^C
      \end{align*}
      de temos (1) e $(E\cap F) \subset F$, segue que diretamente de $\mu(F\E) = 0$,
      que $\mu(F) = \mu(E \cap F)$. Analogamente, temos $\mu(E) = \mu(F \cap E)$, 
      ou seja, $\mu(F) = \mu(E)$.
  \end{proof}
  \item \textbf{Defina $E\sim F$ se $\mu(E \Delta F) = 0$. Mostre que $\sim$ é uma 
    relação de equivalência.}
    \begin{proof}[\textbf{Dem:}]
      Vejamos que vale a reflexiva, simétrica e transitiva. 
     
      \noindent
      (Reflexiva:) Já que para todo $A \in \mcal F$, $A \Delta A = (A\setminus A) \cup (A \setminus A) = \emptyset$, 
      então $\mu(A\Delta A) = 0$, i.e, $A \sim A$;
      
      \noindent 
      (Simétrica:) Suponhamos que, dados $A, B \in \mcal F$, temos $A \sim B$. Ora, como $A \sim B$, 
      $\mu(A\Delta B) = 0$, entrentanto,
      \begin{align*}
        A\Delta B = (A\setminus B) \cup (B \setminus A) = (B \setminus A) \cup (A \setminus B) = B\Delta A,
      \end{align*}
      logo, $\mu(A\Delta B) = \mu(B\Delta A) = 0$, i.e, $B \sim A$.

      \noindent 
      (Transitiva:) Sejam $A, B, C \in \mcal F$, vejamos que se $A\sim B$ e $B \sim C$, então $A \sim C$, ou seja,  
      $\mu (A\Delta C) = 0$. Note que, 
      \begin{align*}
        A\Delta C &= (A\cap C^C) \cup (C \cap A^C) = (A \cap C^C\cap(B \cup B^C)) \cup (C \cap A^C\cap(B \cup B^C))\\
        &=((A\cap C^C\cap B) \cup (A \cap C^C \cap B^C)) \cup ((C\cap A^C \cap B)\cup (C\cap A^C \cap B^C))\\
        &=((A\cap(B \setminus C))\cup(C^C\cap(A\setminus B)))\cup((C\cap(B\setminus A))\cup (A^C \cap (C \setminus B)))\\
        &\subseteq ((B\setminus C)\cup (A\setminus B))\cup ((B\setminus A) \cup (C\setminus B)) = (A\Delta B) \cup (B \Delta C).
      \end{align*}    
      Portanto, $\mu(A\Delta C) \le \mu((A\Delta B) \cup (B\Delta C)) $, $(2)$, ou seja 
      \begin{align*}
        0 \le \mu(A\Delta C) \le \mu((A\Delta B) \cup (B\Delta C)) \le \mu(A\Delta B) + \mu(B \Delta C) = 0 + 0 = 0,
      \end{align*}
      ou seja $\mu(A \Delta C) = 0$, isto é, $A \sim C$, como queríamos.
    \end{proof}
  \item \textbf{Considere $Y = X/\sim $ o conjunto quociente da 
    relação acima. Mostre que a função 
    $d: Y \times Y \rightarrow \mbb R$ definida por 
    $d([E], [F])\coloneq \mu(E \Delta F)$ define uma métrica.}
    \begin{proof}[\textbf{Dem:}]
      Ora, visto que $\sim$ é reflexiva, $\mu(A \Delta A) = 0$ para todo $A \in \mcal F$, ou seja $d([A], [A]) = 0$ para
      todo $[A] \in Y$. Mais ainda, se $[A] \neq [B]$, i.e, $A \not \sim B$, então $\mu(A\Delta B)\neq 0$, ou mais, 
      $d([A], [B]) = \mu(A \Delta B) > 0$. Da transitividade de $\sim$, $d([A], [B]) = \mu(A \Delta B) = 
      \mu(B \Delta A) = d([B], [A])$. Por fim,  de $(2)$, para cada $[A], [B], [C] \in Y$, 
      \begin{align*}
        d([A], [C]) = \mu(A \Delta C) \le \mu(A \Delta B) + \mu(B \Delta C) = d([A], [B]) + d([B], [C]).
      \end{align*}
      Com isso, temos o que queríamos.
    \end{proof}
\end{enumerate}
\biglinend
\stepcounter{dois}
\myexe{\thedois}
\bigline 
\textbf{Sejam $(X, \mcal F, \mu)$ um espaço de medida completo e $(f_n)_{n \in \mbb N}$ uma sequência de funções mensuráveis 
$f_n: X \rightarrow \xbR$. Mostre que se $f_n \rightarrow f$, em quase-$\mu$ todo ponto para uma função 
$f: X \rightarrow \xbR$, então $f$ é mensurável.}
\begin{proof}[\textbf{Dem:}]Como a propriedade vale $\mu$-q.t.p, então, existe $N \in \mcal F$ tal que $\mu(N) = 0$ e 
  vale que $f_n \rightarrow f$ para $x \in N^C$. Considere $A = \{x \in X\st f_n\text{ não converge para }f\}$. Visto que 
  $A \subset N$ e $\mu(N) = 0$, da completude do espaço, temos $A \in \mcal F$. Mais ainda, dado $\alpha \in \mbb R$
  \begin{align*}
    \{x \in A \st f(x) > \alpha\} \subset A \subseteq N
  \end{align*}
  temos novamente da completude do espaço que $\{x \in A \st f(x)>\alpha\} \in \mcal F$. Dito isso, 
  \begin{align*}
    \{x \in X \st f(x)> \alpha \} &= \{x \in A \st f(x)>\alpha\}\cup \{x \in A^C \st f(x) > \alpha\}\\
    &= \{x \in A \st f(x)>\alpha\}\cup \left(\bigcup\limits_{m \in \mbb N} \left[
      \bigcap\limits_{n = m}^{\infty}\{x \in A^C \st f_n(x) > \alpha\} \right]\right)\\
    &= \{x \in A \st f(x)>\alpha\}\cup \left(\bigcup\limits_{m \in \mbb N} \left[
      \bigcap\limits_{n = m}^{\infty}A^C\cap\{x \in X \st f_n(x) > \alpha\} \right]\right)\\
    &= \{x \in A \st f(x)>\alpha\}\cup \left(A^C \cap \bigcup\limits_{m \in \mbb N} \left[
      \bigcap\limits_{n = m}^{\infty}\{x \in X \st f_n(x) > \alpha\} \right]\right). 
  \end{align*}
  Visto que $A \in \mcal F$, temos $A^C \in \mcal F$. Mais ainda, $\lim \inf \{x \in X; f_n(x)> \alpha\} \in \mcal F$, 
  ou seja, $f$ é mensurável.
\end{proof}
\biglinend
\newpage
\stepcounter{dois}
\myexe{\thedois}
\bigline 
\textbf{Seja $(X, \mcal F, \mu)$ um espaço de medida e considere $\mcal Z\coloneq \{N \in \mcal F\st \mu(N) = 0\}$. 
Mostre que se $N \in \mcal Z$ e $F \in \mcal F$, então $N \cap F \in \mcal Z$. Além disso, mostre que se 
$(N_k)_{k \in \mbb N}$ é uma sequência em $\mcal Z$, então $\bigcup_{k \in \mbb N} N_k \in \mcal Z$.}
\begin{proof}[\textbf{Dem:}]
 Já que $N \cap F \subset N$, 
  \begin{align*}
    \mu(N \cap F) \le \mu(N) = 0 \implies N\cap F \in \mcal Z. 
  \end{align*}
  Vejamos que $\bigcup N_k \in \mcal Z$. Ora, da subaditividade, 
  \begin{align*}
    \mu\left(\bigcup\limits_{k \in \mbb N} N_k\right) \le \sum_{k \in \mbb N} \mu(N_k)
  \end{align*}
  Já que $(N_k)$ é uma sequência em $\mcal Z$, $\mu(N_k) = 0$ para todo $k \in \mbb N$. Ou seja, 
  $\bigcup N_k \in \mcal Z$, como queríamos.
\end{proof}
\biglinend 
\stepcounter{dois}
\myexe{\thedois}
\bigline
\textbf{Seja $\overline{\mcal F} \coloneq \{E\cup Z\st E \in \mcal F, Z \subset N\text{ para algum }N\in \mcal Z\}$.
Mostre que $\overline{\mcal F}$ é uma $\sigma$-álgebra, chamada de complemento de $\mcal F$ com respeito a $\mu$}.
\begin{proof}[\textbf{Dem:}]
Visto que $\mu(\emptyset) = 0$, temos $\emptyset \in \mcal Z$. E, como $\emptyset \in \mcal F$, segue que 
  $\emptyset \cup \emptyset = \emptyset \in \overline{\mcal F}$. Mais ainda, visto que $X \in \mcal F$ e, 
  para todo $Z \subseteq N \in \mcal Z \subset \mcal F$, temos $X \cup Z = X$, então $X \in \overline{\mcal F}$. 
  Daqui, vejamos que se $A \in \overline{\mcal F}$, então $A^C \in \overline{\mcal F}$.
  
  \noindent 
  Note que, se $A \in \overline{\mcal F}$, então existem $E \in
  \mcal F$ e $Z \subset N\in\mcal Z$
  tais que $A = E \cup Z$. Portanto, 
  \begin{align*}
    A^C &= (E \cup Z)^C = E^C \cap Z^C = E^C\cap Z^C \cap (N \cup N^C)\\
    &= (E^C \cap Z^C \cap N) \cup (E^C \cap Z^C \cap N^C).
  \end{align*}
  Já que $Z \subseteq N$, temos $N^C \subseteq Z^C$, ou seja, 
  \begin{align*}
    A^C &= (E^C \cap Z^C \cap N) \cup (E^C \cap N^C).
  \end{align*}
  Daqui, note que $\tilde Z = E^C \cap Z^C \cap N \subset N$ e $\tilde A = E^C \cap N^C \in 
  \mcal F$, isto é, $A^C = \tilde A \cup \tilde Z \in \overline{\mcal F}$. 

  \noindent
  Vejamos que dada uma sequência $\{A_K\}$ em $\overline{\mcal F}$,
  temos $\bigcup A_k \in \overline{\mcal F}$. Já que, para todo 
  $k \in \mbb N$, temos $A_k \in \overline{\mcal F}$, então, existe $E_k \in \mcal F$ e 
  $Z_k \subset N_k$ para algum $N_k \in \mcal Z$. Com isso, 
  \begin{align*}
    \bigcup_{k \in \mbb N} A_k = \bigcup_{k \in \mbb N} (E_k \cup Z_k) = 
    \left(\bigcup_{k \in \mbb N} E_k \right)\cup \left(\bigcup_{k \in \mbb N} Z_k\right).
  \end{align*}
  Daqui, note que $\bigcup E_k \in \mcal F$ e $\bigcup Z_k \subseteq 
  \bigcup N_k \in \mcal Z$, portanto $\bigcup A_k \in \overline{\mcal F}$.
\end{proof}
\biglinend 
\stepcounter{dois}
\myexe{\thedois}
\bigline
\textbf{Considere a aplicação $\overline \mu: \overline{\mcal F} \rightarrow \xbR$ dada por 
\begin{align*}
  \overline \mu (E\cup Z) \coloneq \mu(E),
\end{align*}
em que $E\in \mcal F$ e $Z \subset N$, para algum $N \in \mcal Z$. Mostre que:} 
\begin{enumerate}[label=(\alph*), wide, labelindent=0pt]
  \item \textbf{$\overline \mu$ está bem definida, isto é, 
    se $E_1 \cup Z_1 = E_2 \cup Z_2$ então $\overline \mu (E_1 \cup Z_1) = 
    \overline \mu (E_2 \cup Z_2)$}.
    \begin{proof}[\textbf{Dem:}]Ora, dados $N_1, N_2 \in \mcal Z$ tais que $Z_1 \subset N_1$ 
      e $Z_2 \subset N_2$, temos,
      \begin{align*}
        \overline{\mu}(E_1\cup Z_1) &= \mu(E_1) \le \mu(E_1 \cup N_1)
        \le \mu((E_1 \cup Z_1) \cup N_1)\\
        &=\mu((E_2 \cup Z_2) \cup N_1) = \mu(E_2 \cup (N_1 \cup N_2)).
      \end{align*}
      Visto que $N_1 \cup N_2 \in \mcal Z$, temos,
      \begin{align*}
        &\le\mu(E_2) + \mu(N_1 \cup N_2) = \mu(E_2) = \overline{\mu}(E_2 \cup Z_2), 
      \end{align*}
      analogamente, segue o outro lado da desigualdade. Com isso, temos o que queríamos.
    \end{proof}
  \item \textbf{$\overline \mu$ é uma medida em $(X, \overline{\mcal F})$}. 
    \begin{proof}[\textbf{Dem:}]
      Note que $\overline\mu(\emptyset) = \mu(\emptyset) = 0$ e para todo $A \in 
      \overline{\mcal F}$, existe $E \in \mcal F$ e $Z \subset N \in \mcal Z$ tal que 
      $A = E \cup Z$, 
      segue que $\overline \mu(A) = \mu(E) \ge 0$. Com isso, basta verificar a 
      $\sigma$-aditividade para concluír que $\overline \mu$ é uma medida em $(X, \mcal F)$.
      Para isso, considere ${A_k}$ uma sequência em $\overline{\mcal F}$ 
      disjunto dois a dois. Então, para 
      cada $k \in \mbb N$, existe $E_k \in \mcal F$ e $Z_k \subset N_k \in \mcal Z$ 
      para algum $N_k$, desse modo, já que $\bigcup N_k \in \mcal Z$ e 
      $\bigcup Z_k \subseteq \bigcup N_k$, 

      \begin{align*}
        \overline \mu\left(\bigcup_{k \in \mbb N} A_k \right) &=
        \overline\mu\left(\left[\bigcup_{k \in \mbb N} E_k \right] \cup 
        \left[\bigcup_{k \in \mbb N} Z_k \right]\right)\\ 
        &= \mu\left(\bigcup_{k \in \mbb N} E_k\right) = 
        \sum_{k \in \mbb N} \mu(E_k)\\
        &= \sum_{k \in \mbb N} \overline \mu(E_k \cup Z_k)
        =\sum_{k \in \mbb N} \overline \mu(A_k),
      \end{align*}
      como queríamos.
    \end{proof}
  \item \textbf{o espaço de medida $(X, \overline{\mcal F}, \overline \mu)$ é completo.}
    \begin{proof}[\textbf{Dem:}] Considere $A \in \overline{\mcal F}$ satisfazendo 
      $\overline{\mu}(A) = 0$. Basta observar que $A \in \mcal Z$ e 
      todo subconjunto $Z \subset A$
      pode ser escrito como $\emptyset \cup Z \in \overline{\mcal F}$, já que 
      $\emptyset \in \mcal F$. 
            
    \end{proof}
  \item \textbf{$\overline \mu$ coincide com $\mu$ em $\mcal F$}
    \begin{proof}[\textbf{Dem:}]Ora, dado $A \in \mcal F$, sabendo que 
      $\emptyset \subset \emptyset \in \mcal Z$, temos $A \in \overline {\mcal F}$,
      desse modo, 
      \begin{align*}
        \overline \mu(A) = \overline \mu(A \cup \emptyset) = \mu(A).
      \end{align*}
      Portanto, $\overline \mu|_{\mcal F} = \mu$.
    \end{proof}
\end{enumerate}
\biglinend 
\stepcounter{dois}
\myexe{\thedois}
\bigline
\textbf{Utilizando propriedades da medida exterior $m^*$ em $\mbb R$, mostre 
que $[0, 1]$ é não-enumerável. Conclua que $\mbb R$ é não-enumerável.
}
\begin{proof}[\textbf{Dem:}] A fim de contradição, suponhamos que $[0, 1]$ é 
  enumerável. Sabemos que todo conjunto enumerável tem medida exterior zero,
  portanto $m^*([0, 1]) = 0$. Entretanto, $[0, 1]$ é também um intervalo,
  sendo assim, $m^*([0, 1]) = 1-0 = 1$, uma contradição. De maneira análoga
  ao apresentado, a fim de contradição, suponha $\mbb R$ enumerável, 
  ou seja, $m^*(\mbb R) = 0$. 
  Entretanto, $[0, 1] \subset \mbb R$, ou seja, $m^*(\mbb R) \ge 1$, 
  uma contradição.
  
\end{proof}
\biglinend 
\stepcounter{dois}
\myexe{\thedois}
\bigline
\textbf{Utilizando propriedades da medida exterior, prove que se $m^*(A) = 0$, 
então $m^*(A\cup B) = m^*(B)$.}
\begin{proof}[\textbf{Dem:}] Ora, dados $A, B \subset \mbb R$, temos $B \subset (A \cup B)$,
  portanto, 
  \begin{align*}
    m^*(B) \le m^*(A\cup B) \le m^*(A) + m^*(B).
  \end{align*}
  Já que $m^*(A) = 0$, temos o que queríamos.
\end{proof}
\biglinend
\newpage
\stepcounter{dois}
\myexe{\thedois}
\bigline
\textbf{Verifique que os seguintes itens são equivalentes:
\begin{enumerate}
  \item $E$ é mensurável. 
  \item Para todo $\epsilon > 0$, existe um fechado $F$, com 
    $F \subseteq E$, tal que $m^*(E\setminus F) < \epsilon$. 
  \item Existe um conjunto $F_{\sigma}$, denotado por $H$, 
    com $H \subseteq E$, tal que $m^*(E\setminus H) = 0$
\end{enumerate}}
\bigline
\begin{proof}[\textbf{Dem:}] Considere $A \in \xbR$ e seja $E = A^C$. 

\noindent
  $((1) \implies (2))$ Se $E$ é mensurável, então $A$ também o é. Desse modo, 
  para todo $\epsilon > 0$, existe um aberto $\mcal O$, com $A \subseteq \mcal O$, 
  tal que $m^*(\mcal O \setminus A) < \epsilon$. Agora, considere $F = \mcal O^C$. 
  Como $F \subseteq A^C = E$, e, 
  \begin{align*}
    \mcal O \setminus A = \mcal O \cap A^C = \mcal O \cap E = E\cap F^C = E\setminus F
  \end{align*}
  temos $m^*(E\setminus F) < \epsilon$, como queríamos.

  \noindent
  $((2) \implies (3))$ Suponhamos que para todo $\epsilon > 0$ existe um fechado $F$,
  com $F \subseteq E$, tal que $m^*(E\setminus F) < \epsilon$. Então, para cada $k \in \mbb N$,
  tome $F_k$ fechados com $F_k \subseteq E$, tais que $m^*(E\setminus F_k)< 
  \frac{1}{k}$. Daqui, tomando $H = \bigcup F_k$, note que $H \subseteq E$, pois caso contrário
  existiria $F_k \not \subseteq E$, uma contradição. Já que $E\setminus H 
  \subseteq E\setminus F_k$, para todo $k \in \mbb N$
  \begin{align*}
    m^*(E\setminus H) < m^*(E\setminus F_k) < \frac{1}{k}. 
  \end{align*}
  sendo assim, $m^*(E\setminus H) = 0$, como queríamos.
  
  \noindent
  $((3) \implies (1))$ Suponhamos que existe um conjunto $F_\sigma$, denotado 
  por $H$, com $H \subseteq E$, tal que $m^*(E\setminus H) = 0$. Visto que $H$ pode 
  ser escrito por uniões de intervalos, e, todo intervalo é mensurável, temos 
  que $H$ é mensurável. Mais ainda, visto que $m^*(E \setminus H) = 0$, 
  temos $E\setminus H$ mensurável. Daqui, vejamos que $E$ é mensurável. Para isso, 
  note que, 
  \begin{align*}
    (E\setminus H)^C\cap H^C = (E \cap H^C)^C \cup H^C = (E^C \cup H) \cap H^C
    = E^C,
  \end{align*}
  visto que $H\subseteq E \implies E^C \subseteq H^C$. Com isso, $E^C$ é mensurável,
  portanto $E$ é mensurável.
\end{proof}
\biglinend
\newpage
\stepcounter{dois}
\myexe{\thedois}
\bigline 
\textbf{Mostre que um conjunto $E$ é mensurável se, e somente se, para cada $\epsilon > 0$, 
existe um conjunto fechado $F$ e um conjunto aberto $\mcal O$, tais que 
$F \subseteq E \subseteq \mcal O$ e $m^*(\mcal O\setminus F) < \epsilon$.}
\bigline
\begin{proof}[\textbf{Dem:}]
  \hfill 
  
  \noindent 
  $(\implies)$ Ora, visto que $E$ é mensurável, para todo $\epsilon > 0$, exisem conjuntos 
  $\mcal O$ e $F$, aberto e fechado, respectivamente, tais que 
  $F\subseteq E \subseteq \mcal O$ e $m^*(\mcal O\setminus E) < \frac{\epsilon}{2}$
  e $m^*(E\setminus F) < \frac{\epsilon}{2}$. Daqui, temos
  \begin{align*}
    \epsilon \ge m^*(\mcal O\setminus E) + m^*(E\setminus F) \ge m^*(\mcal O \setminus F), 
  \end{align*}
  visto que, 
  \begin{align*}
    (\mcal O \setminus E) \cup (E \setminus F) &= 
    (\mcal O \cap E^C) \cup (E \cap F^C) = [\mcal O \cup (E \cap F^C)] \cap [\E^C \cup 
    (E \cap F^C)]\\ 
    &= [\mcal O\cup \emptyset] \cap [\xbR \cap F^C]
    = \mcal O \cap F^C = \mcal O \setminus F.
  \end{align*}
  Como queríamos.

  \noindent 
  $(\impliedby)$ Por outro lado, suponhamos que para todo $\epsilon>0$ existe 
  um conjunto fechado $F$ e um conjunto aberto $\mcal O$, tais que 
  $F \subseteq E \subseteq \mcal O$ e $m^*(\mcal O\setminus F) < \epsilon$.
  Visto que $(\mcal O\setminus E) \cap (E\setminus F) = \emptyset$, segue que 
  \begin{align*}
    \epsilon &> m^*(\mcal O\setminus F) = m^*(\mcal O \setminus E) + m^*(E \setminus F)\\
    &\ge m^*(\mcal O\setminus E).
  \end{align*}
  Com isso, temos $E$ mensurável.
\end{proof}
\biglinend
\newpage
\stepcounter{dois}
\myexe{\thedois}
\bigline 
\textbf{Considere a reta real $\mbb R$ munida da $\sigma$-álgebra de Borel 
$\mcal B (\mbb R)$ e seja $\mu_B$ a restrição da medida exterior $m^*$ aos 
conjuntos de Borel. Seja $(\mbb R, \overline{\mcal B (\mbb R)}, \overline{\mu}_b)$, conforme 
construído nos Exercícios $12$ a $14$. O objetivo deste exercício é mostrar que 
o espaço de medida de Lebesgue $(\mbb R, \mcal M, m)$ é precisamente este complemento. 
Prove os seguintes itens:
\begin{enumerate}[label=(\alph*), wide, labelindent=0pt]
  \item Mostre que $\overline{\mcal B(\mbb R)} \subseteq \mcal M$. 
  \begin{proof}[\textbf{Dem:}]
   Por definição, 
    \begin{align*}
      \overline{\mcal B(\mbb R)} = \{E \cup Z\st E\in \mcal B(\mbb R), Z\subset N 
      \text{ para algum } N\in\mcal Z\},
    \end{align*}
    onde $\mcal Z = \{N \in \mcal B(\mbb R)\st \mu_B(N) = 0\}$. Visto que 
    $\mcal B(\mbb R) \subseteq \mcal M$, temos que todo elemento da $\sigma$-álgebra 
    de Borel é mensurável. Daqui, tome $A \in \overline{\mcal B(\mbb R)}$, desse modo
    existem $E \in \mcal B(\mbb R)$ e $Z \subset N \in \mcal Z$ tais que $A = E \cup Z$.
    Já que $\overline{\mcal B(\mbb R)}$ é completo, $\mu_B(Z) = m^*(Z) = 0$, ou seja $Z$ é 
    mensurável. Como $A$ pode ser escrito como a união de elementos de $\mcal M$ 
    e $\mcal M$ é $\sigma$-álgebra, temos $A \in \mcal M$, portanto, 
    $\overline{\mcal B (\mbb R)} \subseteq \mcal M$. 
  \end{proof}
\item Utilizando o exercício anterior, mostre que todo conjunto $E \in \mcal M$ 
  pode ser escrito na forma $E = F\cup Z$, em que $F \in \mcal B(\mbb R)$ e $Z \subset N$ 
    para algum $N \in \mcal B(\mbb R)$ com $\mu_B(N) = 0$. Conclua que 
    $\mcal M = \overline{\mcal B(\mbb R)}$ e que $m(E) = \overline{\mu}_B(E)$ para todo 
    $E \in \mcal M$.
    \begin{proof}[\textbf{Dem:}] Considere $E$ mensurável. Visto que existe $F_\sigma$
      denotado por $H\in \mcal B(\mbb R)$, com $H \subseteq E$, tal que $m^*(E\setminus H) = 0$, segue que 
      $C = E\setminus H$  é mensurável e $E = H \cup C$. 
      Daqui, note que existe $G_\delta$ denotado por $\mcal O$ tal que $\mcal O \supset C$
      e $m(\mcal O\setminus C) = 0$. Com isso, visto que $\mcal O \in \mcal B(\mbb R)$,
      \begin{align*}
        m(\mcal O \setminus C) = m(\mcal O) - m(C) = m(\mcal O) = \mu_B(\mcal O) = 0,
      \end{align*}
      ou seja, $E\in \mcal M$ pode ser escrito na forma $E = H \cup C$ tal que $H \in \mcal B(\mbb R)$
      e $C\subset \mcal O$ tal que $\mcal O \in \mcal B(\mbb R)$ com $\mu_B(\mcal O) = 0$, i.e, 
      $E \in \overline{\mcal B(\mbb R)}$. Por fim, tome $A \in \mcal M$, portanto, existem 
      $E$ e $Z$ satisfazendo as condições do enunciado. Daqui, note que
      \begin{align*}
        m(A) = m(E\cup Z) \le m(E) + m(Z) = m(E) = m^*(E) = \mu_B(E) = \overline\mu_B(E\cup Z) = \overline\mu_B(A).
      \end{align*}
      Por outro lado, 
      \begin{align*}
        \overline\mu_B(A) = \overline\mu_B(E\cup Z) = \mu_B(E) = m^*(E) = m(E) \le m(E\cup Z) = m(A).
      \end{align*} 
      Ou seja, $m = \overline \mu_B$ em $\mcal M$.
    \end{proof}
\end{enumerate}}
\biglinend
\newpage
\stepcounter{dois}
\myexe{\thedois}
\bigline 
\textbf{Mostre que se $E_1$ e $E_2$ são conjuntos Lebesgue mensuráveis, então 
\begin{align*}
  m(E_1 \cup E_2) + m(E_1 \cap E_2) = m(E_1) + m(E_2).
\end{align*}}
\begin{proof}[\textbf{Dem:}] Ora, se $m(E_1 \cap E_2) = +\infty$, segue que $E_1$ 
  e $E_2$ tem medida infinita, portanto vale a igualdade. Suponhamos que $m(E_1 \cap E_2)<+\infty$,
  então, 
  \begin{align*}
    m(E_1\cup E_2) &= m([E_1 \cup E_2]\cap E_1) + m([E_1 \cup E_2]\cap E_1^C)\\
    &= m(E_1) + m(E_2 \cap E_1^C).\\
  \end{align*}
  Já que, 
  \begin{align*}
    E_2 \setminus (E_1 \cap E_2) = E_2 \cap (E_1^C \cup E_2^C) = E_2\cap E_1^C,
  \end{align*}
  segue do fato que $(E_1 \cap E_2) \subset E_2$ que
  \begin{align*}
    m(E_1 \cup E_2) = m(E_1) + m(E_2 \setminus(E_1 \cap E_2)) = m(E_1) + m(E_2) - m(E_1 \cap E_2).
  \end{align*}
  Como $m(E_1 \cap E_2)<+\infty$, temos o que queríamos.
\end{proof}
\biglinend
\newpage
\stepcounter{dois}
\myexe{\thedois}
\bigline 
\textbf{Seja $(E_n)_{n \in \mbb N}$ uma coleção enumerável de conjuntos Lebesgue mensuráveis disjuntos. Prove 
que para todo $A \in \mcal P(\mbb R)$ vale
\begin{align*}
  m^*\left(A \cap \bigcup\limits_{n \in \mbb N} E_n\right) = \sum_{n \in \mbb N}m^*(A\cap E_n).
\end{align*}}
\begin{proof}[\textbf{Dem:}]
  Ora, se $m^*(A \cap \bigcup E_n) = +\infty$, a igualdade segue diretamente da $\sigma$-subaditividade. Suponhamos 
  que $m^*(A \cap \bigcup E_n) < +\infty$. Seja $(S_k)_{k \in \mbb N}$ tal que $S_k = m^*(A\cap \bigcup_{n=1}^k E_n)$,
  então,
  \begin{align*}
    S_k = m^*\left(\bigcup_{n=1}^k A \cap E_n\right) \le m^*\left(\bigcup_{n \in \mbb N} A\cap E_n\right).
  \end{align*}
  Como $(E_n)$ é composto por conjuntos disjuntos, 
  temos que $(A \cap E_n)_{n \in \mbb N}$ é disjunto, portanto, para todo $k \in \mbb N$ temos,
  \begin{align*}
    S_k = \sum_{n = 1}^k m^*(A \cap E_n) \le m^*\left(\bigcup_{n \in \mbb N} A\cap E_n\right).
  \end{align*}
  Já que $S_k$ é não-decrescente e limitada superiormente por $m^*(\bigcup A\cap E_n)$, temos que 
  o limite e limitado pelo mesmo, portanto, com a $\sigma$-subaditividade, segue que
  \begin{align*}
    \sum_{n \in \mbb N} m^*(A\cap E_n) =  m^*\left(\bigcup_{n \in \mbb N} A\cap E_n\right),
  \end{align*}
  como queríamos.
\end{proof}
\biglinend
\newpage
\stepcounter{dois}
\myexe{\thedois}
\bigline 
\textbf{Seja $E \subseteq \mbb R$ um conjunto Lebesgue mensurável com $m(E)>0$. Fixe um número real $L$ tal que 
$0 < L < m(E)$. Defina a função $f: [0, +\infty) \rightarrow \mbb R$ dada por $f(x) = m(E\cap[-x, x])$
\begin{enumerate}[label=(\alph*), wide, labelindent=0pt]
  \item Mostre que $f$ é uma função Lipschitziana.
    \begin{proof}[\textbf{Dem:}]
        Sejam $x, a \in [0, +\infty)$. Sem perda de generalidade, suponha $x \ge a$. 
        Desse modo, 
      \begin{align*}
        |f(x) - f(a)| = |m(E\cap [-x, x]) - m(E\cap [-a, a])|.
      \end{align*}
      Já que $E\cap [-a,a] \subset E\cap [-x, x]$, então, 
      \begin{align*}
        |f(x) - f(a)| &= |m(E\cap [-x, x] \setminus [E\cap [-a, a]])|\\ 
        &= |m([-x, -a] \cup [a, x])| = |m([-x, -a]) + m([a, x])|\\ 
        &= |-a + x +x - a| = 2|x-a|.
      \end{align*}
      Com isso, segue que $f$ é uma função Lipschitziana.
    \end{proof}
  \item Determine os valores de $f(0)$ e $\lim\limits_{x \rightarrow \infty} f(x)$.
    \begin{proof}[\textbf{Solução:}]
      Para determinar o valor de $f(0)$, note que $E\cap[-0, 0] = E\cap\{0\}$ é 
      igual a $\{0\}$ ou $\emptyset$.
      Se $E\cap\{0\} = \{0\}$, temos $f(0) = 0$. Caso $E\cap\{0\} = \emptyset$
      temos $f(0)=0$. Desse modo, $f(0)=0$.\\
      \noindent
      Para determinar o valor de $\lim_{x\rightarrow \infty} f(x)$, note que 
      $E \subseteq \mbb R = (-\infty, \infty)$. Portanto, 
      \begin{align*}
        \lim_{x \rightarrow \infty} f(x) = \lim_{x \rightarrow \infty} m(E \cap [-x, x])
        = m(E\cap(-\infty, \infty)) = m(E).
      \end{align*}
      Com isso, temos o que queríamos.
    \end{proof}
    \item Utilize o Teorema do Valor Intermediário para concluir que existe um subconjunto
      mensurável $B_L \subseteq E$ tal que $m(B_L) = L$.
      \begin{proof}[\textbf{Dem:}]
        Já que $f$ é Lipschitziana, temos $f$ uniformemente contínua.
        Desse modo, do Teorema do Valor 
        Intermediário, visto que $L \in (0, m(E)) = (f(0), \lim f(x))$ e $[0, +\infty)$
        é fechado, segue que existe $c \in [0, +\infty)$ tal que $f(c) = L$. Portanto, 
        \begin{align*}
          L = f(c) = m(E\cap[-c, c]),
        \end{align*}
        isto é, $B_L = E\cap [-c, c]$.
      \end{proof}
\end{enumerate}
}

